\documentclass[10pt,a4paper]{amsart}
\usepackage[margin=19mm]{geometry}
\usepackage{amsmath,amssymb,mathtools}
\usepackage[scr=rsfs,cal=euler]{mathalfa}
\usepackage{xfrac}
\usepackage[unicode,hidelinks,bookmarks=false]{hyperref}
\newcommand{\ie}{i.e.\ }
\newcommand{\cf}{cf.\ }
\newcommand{\resp}{respectively\ }
\newcommand{\Subsection}{\subsection}
\providecommand{\nobreakdash}{\nobreak\hbox{-}}
\setlength{\emergencystretch}{2em}
\multlinegap0pt
\usepackage{bm}
\usepackage{bbm}
\usepackage{dsfont}
\usepackage{xspace}
\usepackage{cancel}
\usepackage{mathtools}
\usepackage[shortlabels]{enumitem}
\usepackage{amsthm}
\makeatletter
\@ifundefined{theorem}{\newtheorem{theorem}{Theorem}}{}
\makeatother
\DeclareMathOperator{\dom}{dom}
\DeclareMathOperator{\hht}{ht}
\DeclareMathOperator{\ran}{ran}
\newcommand{\Q}{\mathbb{Q}}
\newcommand{\es}{\emptyset}
\newcommand{\res}{\restriction}
\newcommand{\seb}{\subseteq}
\newcommand{\tup}[1]{\langle#1\rangle}
\renewcommand{\P}{\mathbb{P}}
\title[New consequences of PFA($T^*$)]{New consequences of PFA($T^*$)}
\author{Carlos Mart\'inez-Ranero, Lucas Polymeris}
\date{Source: arXiv:2510.22807v2 (CC BY 4.0)}
\begin{document}
\maketitle

\noindent\emph{Source and licence.} New consequences of PFA($T^*$). Version arXiv:2510.22807v2 (CC BY 4.0), distributed under the Creative Commons Attribution 4.0 International licence, as recorded in the arXiv licence metadata for this exact version and shown on the abstract page https://arxiv.org/abs/2510.22807v2. Full rights notice: licenses/arxiv-2510.22807-CC-BY-4.0.txt.\par\smallskip

\noindent\textit{Editorial scope.} Counted unit: the Theorem whose source label is \texttt{thm:diamond-S} in the article (source0.tex lines 633--644, source label \texttt{thm:diamond-S}), delivered together with the article's own complete original proof of it (source0.tex lines 645--781, one \texttt{proof} environment of 137 lines). No mathematics is written, repaired or re-derived: statement and proof are the article's own text line for line. Required context retained uncounted: none. Every definition, display and label of those lines is retained unchanged. Omitted from this package and not redistributed: 1--632, 782--879. Nothing inside a retained excerpt is omitted or altered: every retained body is delivered exactly as the article's lines are written, so no omission receipt of any kind accompanies this package. Citations: the retained excerpts carry no citation command at all, so no bibliography entry of the article is transcribed and no reference is redistributed. Reference adaptation: none was needed and none was made; every reference inside the delivered unit resolves inside it, so no unresolved reference or citation is produced. Printed slips kept literal: no printed word, symbol, hypothesis or number of the counted statement or of its proof was altered. Layout and class: the article's own class is \texttt{amsart} with packages including lmodern, fontenc, inputenc, amsmath, amssymb, amsthm, enumerate, bbm, multicol, amsfonts, hyperref; the wrapper is the lane's pinned \texttt{amsart} editorial wrapper at 10pt on A4, which restates the theorem-like environments the retained text needs and adds the article's own macro definitions that the retained text needs (\texttt{\textbackslash P}, \texttt{\textbackslash Q}, \texttt{\textbackslash dom}, \texttt{\textbackslash es}, \texttt{\textbackslash hht}, \texttt{\textbackslash ran}, \texttt{\textbackslash res}, \texttt{\textbackslash seb}, \texttt{\textbackslash tup}), with extra wrapper packages (bm, bbm, dsfont, xspace, cancel, mathtools, enumitem). The delivered numbering is the unit's own. Assets: no retained span contains a figure environment, a graphics-inclusion command, a PDF-page inclusion, a table or a TikZ picture; no image file is copied into this package, the package declares no asset dependency and no image file is redistributed with it. Licence: version arXiv:2510.22807v2 of this article is distributed under the Creative Commons Attribution 4.0 International licence, as recorded in the arXiv licence metadata for this exact version and shown on the abstract page https://arxiv.org/abs/2510.22807v2. A full rights notice is delivered at licenses/arxiv-2510.22807-CC-BY-4.0.txt. Characters: every retained excerpt is pure ASCII, so the delivered engine, XeLaTeX with its default Latin Modern text font, reports no missing character.
\begin{theorem}\label{thm:diamond-S}
  Assume $\diamondsuit$. For every $1 \le d < \omega$,
  there is a normal Suslin tree
  $S$ such that:
  \begin{itemize}
    \item All its derived trees of dimension at most $d$ are
    Suslin.
    \item All its derived trees of dimension greater than $d$
    are special.
  \end{itemize}

\end{theorem}

\begin{proof}
  Fix $d \ge 1$. First observe  that it is enough to make sure that all derived
  trees of dimension exactly $d$ are Suslin, and all derived trees
  of dimension exactly $d+1$ are special.

  Let us fix some notation. For $1 \le n < \omega$, let
  $S^{(n)}$ denote the $\bar{x} \in \bigotimes_{i < n}S$ that
  are injective, i.e., without repetitions. Similarly
  define $S^{(n)}_{\alpha}$. We also understand that
  if $\bar{x} \in S^{(n)}$, then $\bar{x} = (x_{0},\dots,x_{n-1})$.
  We order $S^{(n)}$ with the partial order induced
  from $\bigotimes_{i<_{n}}S$, which we will denote simply
  by $<_{S}$.

  Fix $\tup{A_{\delta} : \delta < \omega_{1}}$ a $\diamondsuit$-sequence
  for subsets of $\omega_{1}^{d}$.
  We will construct $S$ by recursion on the levels, making
  sure that for each $\alpha < \omega_{1}$,
  $S\res \alpha$ is a countable ordinal as a set.
  Together with $S$ we will define
  $f : \bigotimes_{i < d+1}S \to \Q_{\ge 0}$ satisfying the following:
  \begin{enumerate}[(i)]
    \item $f(\bar{x}) = 0$ iff $\bar{x}$ is not injective, i.e.,
    $\bar{x} \notin S^{(d+1)}$.
    \item For all $\bar{x},\bar{y} \in S^{(d+1)}$,
    if $\bar{x} <_{S} \bar{y}$ then $f(\bar{x}) < f(\bar{y})$.
    \item Suppose $\alpha < \beta$, and
    $\bar{x} \in S_{\alpha}^{(n)}$ with $n \ge d$. Then for all
    $I \in [n]^{d}$ and $\sigma : I \to S_{\beta}$ such that
    for all $i \in I$, $x_{i} <_{S} \sigma(i)$, and for
    all $\varepsilon \in \Q^{+}$, there is
    $\bar{y} \in S_{\beta}^{(n)}$ such that $\bar{x} <_{S} \bar{y}$,
    $y_{i} = \sigma(i)$ for $i \in I$,
    and for all $i_{0},\dots,i_{d} < n$ without
    repetitions,
    $f(x_{i_{0}},\dots,x_{i_{d}}) < f(y_{i_{0}},\dots,y_{i_{d}})
    < f(x_{i_{0}},\dots,x_{i_{d}}) + \varepsilon$.
  \end{enumerate}
  Note that property (ii) will ensure that all derived trees of
  dimension $d+1$ are special.

  Let $S_{0} := \{0\}$ and given
  $S_{\alpha}$, construct $S_{\alpha+1}$ by adding $\aleph_{0}$
  immediate successors to each node in $S_{\alpha}$ so that
  (i)--(iii) are satisfied. Now
  suppose $S \res \delta$
  has been defined for some countable limit ordinal satisfying (i)--(iii),
  and all the normality requirements. We will construct
  a countable set $B$ of cofinals branches through $S \res \delta$
  satisfying the following:
  \begin{itemize}
    \item[(B1)] For every $x \in S \res \alpha$ there is $b \in B$ such
    that $x \in B$.
    \item[(B2)] If it happens that $A_{\delta}$ is a maximal
    antichain of $S(x_{0}) \otimes \cdots \otimes S(x_{d-1})$ for
    some $\bar{x} \in (S\res \delta)^{(d)}$, then for every
    $b_{0},\dots,b_{n-1} \in B$ such that for all $i < d$,
    $x_{i} \in b_{i}$, $\bigotimes_{i < d}b_{i} \cap A_{\delta} \neq \es$.
  \end{itemize}
  We will define $S_{\delta}:= \{\tilde{b} : b \in B\}$ where for every
  $b \in B$, $\tilde{b}$ is a node above each element of $b$.
  This will ensure that the resulting tree is normal,
  and has every derived tree of dimension $d$ Suslin. Thus,
  it is enough to see how to construct $B$ satisfying (B1) and (B2),
  and how to extend $f$ on $S_{\delta}$ satisfying (i)--(iii).

  Define $\P_{\delta}$ to be the poset given by the
  $(\bar{a},h)$ such that for some $n \in \omega$,
  \begin{itemize}
    \item $\bar{a} \in \bigotimes_{i < n}S$. We call $n$ its \emph{length}
    and denote it by $|\bar{a}|$.
    Observe that we do not require $\bar{a}$ to be injective.
    \item $h : n^{d} \to \Q^{+}$ is such that for all
    $(i_{0},\dots,i_{d}) \in \dom(h)$,
    $f(a_{i_{0}},\dots,a_{i_{d}}) < h(i_{0},\dots,i_{d})$.
  \end{itemize}
  The order is given by $(\bar{b},g) \le (\bar{a},h)$ if
  $|\bar{b}| \ge |\bar{a}|$, $\hht(\bar{b}) \ge \hht(\bar{a})$,
  for all $i < |a|$, $b_{i} \ge_{S} a_{i}$, and
  $g \supseteq h$.

  Let $N \prec H(\omega_{2})$ be countable and such that
  $\delta,S\res\delta,A_{\delta},\P_{\delta} \in N$. Let
  $G$ be a $\P_{\delta}$-generic filter for $N$. For each
  $i < \omega$ let $b_{i} := \{a_{i} : (\bar{a},h) \in G, |\bar{a}| > i\}$.

  We claim that for each $i < \omega$, $b_{i}$ is a cofinal branch
  through $S \res \delta$, and moreover, if
  $B := \{b_{i} : i< \omega\}$, then  we have that (B1) and (B2) hold.
  To see that each $b_{i}$ is a cofinal branch,
  note that for each $i < \omega$
  and $\alpha < \delta$, the set of $(\bar{a},h) \in \P_{\delta}$ such that
  $|\bar{a}| > i$ and $\hht(\bar{a}) \ge \alpha$, is dense and it
  is definable from parameters in $N$, thus it is in $N$, so
  $G$ meets it. Since the  verification that the  clauses  (B1) and (B2) hold, are similar.
  We will only present the proof that (B2) holds. Suppose that $A_{\delta}$ is a maximal
  antichain of $\bigotimes_{i < d}S(x_{i})$, where
  $\bar{x} \in (S\res \delta)^{(n)}$. For each $x_{i}$, let
  $b_{k(i)}$ be such that $x_{i} \in b_{k(i)}$. Observe that
  $k$ is injective since $\bar{x}$ is injective. Now consider the
  set of $(\bar{a},h) \in \P_{\delta}$ such that
  $|\bar{a}| > \sup_{i < d}k(i)$,
  and either $(a_{k(0)},\dots,a_{k(d-1)}) \in A_{\delta}$ or
  it does not extend $\bar{x}$. Then this set is dense by (iii) in
  $S \res \delta$,
  and it is in $N$, thus $G$ meets it. This easily implies what we want.

  It remains to define $f$ on $S_{\delta}$. Call
  $\tilde{b}_{i}$ the node in $S_{\delta}$ above $b_{i}$.
  For every injective $(i_{0},\dots,i_{d})$ we define
  $f(\tilde{b}_{i_{0}},\dots,\tilde{b}_{i_{d}}) := h(i_{0},\dots,i_{d})$
  where $(\bar{a},h) \in G$ is any such that $|\bar{a}| > \sup_{i \le d}i_{d}$.
  This is well defined because $G$ is a filter. Of course if $(i_{0},\dots,i_{d})$
  is not injective we let $f(\tilde{b}_{i_{0}},\dots,\tilde{b}_{i_{d}}) = 0$.
  It follows
  from the definition of $\P_{\delta}$ that this will satisfy
  (i) and (ii). We need to show (iii).

  Pick $\bar{x} \in (S\res \delta)^{n}$. We may assume $n \ge d+1$,
  since otherwise (iii) holds vacuously. Pick $I \in [n]^{d}$,
  $\sigma : I \to \omega$ and $\varepsilon \in \Q^{+}$.
  Consider $D \seb \P_{\delta}$ the set of $(\bar{a},h)$ satisfying that
  \begin{itemize}
    \item $\ran(\sigma) \seb |\bar{a}|$, $\hht(\bar{a}) > \hht(\bar{x})$ and
    $\bar{a}$ is injective.
    \item If it is the case that for all $i \in I$, $x_{i} <_{S} \sigma(i)$,
    then there is $\sigma' : n \to \omega$ such that:
    \begin{itemize}
      \item $\ran(\sigma') \seb |\bar{a}|$ and $\sigma' \res I = \sigma$.
      \item For all injective $(i_{0},\dots,i_{d}) \in n^{d+1}$,
  $h(a_{\sigma'(i_{0})},\dots,a_{\sigma'(i_{d})}) < f(x_{i_{0}},\dots,x_{i_{d}}) + \varepsilon$.
    \end{itemize}
  \end{itemize}
  An application if (iii) on $S \res \delta$ shows that $D$ is dense.
  This easily implies that (iii) holds on $S \res (\delta +1)$ so
  we are done.
\end{proof}

\end{document}
